\documentclass[11pt]{letter}
\usepackage[margin=1in]{geometry}
\usepackage{hyperref}

\signature{Michael Whiteside \\ Independent Researcher \\ ORCID: 0009-0000-0643-5791}
\address{mickwh@msn.com \\ February 2026}

\begin{document}

\begin{letter}{Editor-in-Chief \\ [Journal Name] \\ [Journal Address]}

\opening{Dear Editor,}

We are pleased to submit our manuscript entitled ``\textbf{Tissue-Specific Circadian Gating Signatures Reveal Distinct Cancer Vulnerability Profiles Across Mammalian Organs}'' for consideration for publication in [Journal Name].

\textbf{Summary of Findings:}

Using a novel analytical framework (PAR(2) - Phase-gated Autoregressive analysis), we systematically characterized circadian gating relationships between clock genes and cancer-relevant targets across 12 mouse tissues. Our key findings include:

\begin{enumerate}
    \item \textbf{Tissue-specific gating signatures:} Each tissue deploys the circadian clock to protect distinct pathways---liver gates DNA damage sensing (ATM) and cell cycle checkpoints (Wee1), heart exclusively gates Hippo/YAP growth control (Tead1, Yap1), kidney gates Myc proliferation, and brain regions gate metabolism (Sirt1) and mitosis (Cdk1).
    
    \item \textbf{Cancer vulnerability model:} These patterns explain tissue-specific cancer susceptibilities associated with circadian disruption---shift workers may be at elevated risk for hepatocellular carcinoma (loss of ATM gating), renal cell carcinoma (loss of Myc gating), and brain tumors (loss of Cdk1 gating).
    
    \item \textbf{Chronotherapeutic implications:} Our findings identify potential targets for time-of-day-optimized cancer therapy.
\end{enumerate}

\textbf{Significance:}

While the connection between circadian disruption and cancer is well-established epidemiologically, the molecular basis for tissue-specific vulnerabilities has remained unclear. Our work provides the first systematic comparison of circadian gating across multiple tissues and establishes a ``vulnerability protection'' model that explains why different organs face different cancer risks from clock disruption.

\textbf{Technical Innovation:}

We developed the PAR(2) Discovery Engine, an open-source software tool that enables researchers to test circadian gating hypotheses in their own datasets. This resource will facilitate further investigation of clock-cancer relationships.

\textbf{Data Availability:}

All source data (GSE54650) are publicly available from NCBI GEO. Our analysis software and processed datasets are available at [repository URL].

This manuscript has not been published previously and is not under consideration elsewhere. All authors have approved the manuscript and agree with its submission to [Journal Name].

We suggest the following potential reviewers with expertise in circadian biology and cancer:
\begin{itemize}
    \item [Reviewer 1 Name, Institution, Email]
    \item [Reviewer 2 Name, Institution, Email]
    \item [Reviewer 3 Name, Institution, Email]
\end{itemize}

Thank you for considering our manuscript. We look forward to your response.

\closing{Sincerely,}

\end{letter}
\end{document}
